\chapter{The Shell and the Graphical Interface}
\label{ch:shell-gui}

\section{The shell}

\fn{./morendo -shell} starts the interactive shell on the calling
thread. It reads a line at a time with JLine, so the usual line editing
keys work, Up and Down walk the history, and the history is kept in
\file{\textasciitilde/.morendo\_history} between sessions. When the input
is a pipe rather than a terminal the shell falls back to a plain reader,
which is how the transcripts in this manual were produced:

\begin{lstlisting}[style=shell]
$ ./morendo -shell <<'EOF'
(batch samples/only/only_1.clp)
(fire)
(exit)
EOF
\end{lstlisting}

Input is collected until the parentheses balance, ignoring parentheses
inside strings and after \fn{;;}; the prompt changes to \fn{...} while an
expression is incomplete. Each complete expression is sent to the engine
and everything the engine prints in response, followed by the
expression's value, is written to the terminal. Errors print the
exception with a short stack trace and the shell continues. Ctrl-C
discards a half-typed expression; Ctrl-D, the end of piped input, or
\fn{(exit)} ends the session. Nothing calls \fn{System.exit}: \fn{exit}
closes the engine, and the shell stops when it notices.

The shell does not talk to the engine directly. It opens a
\emph{channel} on the engine's message router and posts each expression
to the router's command thread, which parses and executes it and posts
the output and the result back. The GUI uses the same mechanism, and a
Java program can too (Chapter~\ref{ch:embedding}).

\subsection{Useful commands}

Beyond the constructs of Chapter~\ref{ch:language}, these are the shell
commands to know. All are ordinary functions and can be used in rule
files as well.

\begin{center}
\begin{tabularx}{\textwidth}{lL}
\toprule
\fn{(batch file)} & load a rule file; paths are relative to the working directory, or URLs, or \fn{classpath:} resources \\
\fn{(load-facts file)} & load a data file of facts \\
\fn{(facts)}, \fn{(agenda)}, \fn{(rules)} & list working memory, the agenda and the rules \\
\fn{(templates)}, \fn{(list-deffacts)}, \fn{(modules)}, \fn{(cubes)} & list the other constructs \\
\fn{(functions [group])} & list the functions, all or one group \\
\fn{(usage name)} & a function's usage string \\
\fn{(ppdefrule name)}, \fn{(ppdeffacts name)} & pretty-print a rule or a deffacts \\
\fn{(ppdeftemplate name)}, \fn{(ppdefcube name)} & pretty-print a template or a cube \\
\fn{(watch facts|rules|activations|all)} & trace engine activity; \fn{unwatch} stops it \\
\fn{(profile ...)}, \fn{(print-profile)} & time the engine (Chapter~\ref{ch:analysis}) \\
\fn{(spool name file)}, \fn{(spool off name)} & copy everything printed to a file \\
\fn{(fire [n])} & run the agenda \\
\fn{(reset)} & retract every fact, assert the deffacts and the objects again \\
\fn{(reset-facts)}, \fn{(reset-objects)} & retract and re-assert the current facts or objects \\
\fn{(clear)} & empty the engine \\
\fn{(version)}, \fn{(mem-used)}, \fn{(gc)} & housekeeping \\
\fn{(exit)} & leave \\
\bottomrule
\end{tabularx}
\end{center}

\subsection{Output}

Everything a rule prints with \fn{printout t} goes to the shell's
channel; output to any other router name goes to the file opened under
that name with \fn{open}, and to nothing else
(Section~\ref{sec:routers}). \fn{readline t} and \fn{read t} inside a
command take the next line typed at the shell. With \fn{(watch facts)} on, every assertion is echoed as
\fn{==> f-n (...)} and every retraction as \fn{<== f-n (...)}; with
\fn{(watch rules)} each firing prints
\fn{==> fire: Activation: rule, id-n, ...} naming the facts that made the
match; \fn{(watch activations)} prints \fn{=> Activation: ...} as
activations reach the agenda. Log messages from the engine's Log4j
loggers go to standard error, at level WARN unless
\fn{-Dmorendo.log.level=DEBUG} is passed to the JVM; add
\fn{-Dlog4j2.configurationFile=classpath:log4j2-file.xml} to also write
\file{logs/morendo.log}. To pass JVM options, edit the \fn{morendo}
script or run \fn{java -cp "libs/*" org.morendo.Morendo -shell} yourself.

\section{The graphical interface}

\fn{./morendo -gui} opens a Swing window with seven tabs. With
\fn{-gui} alone, closing the window closes the engine and exits the JVM;
with \fn{-gui -shell} the window and the shell share one engine and the
JVM exits when both are gone. Window size and position are remembered
between runs.

\begin{description}[style=nextline]
\item[Shell] A console with the same behaviour as the terminal shell:
  the \fn{Morendo>} prompt, a hundred lines of history on Up and Down,
  copy and paste, and a Clear button. Font and colours are set on the
  Settings tab.
\item[Facts] A sortable table of every fact with its id. Reload refreshes
  it; Assert Fact opens an editor that lists the templates and builds an
  \fn{assert} from the values you enter; right-clicking a row retracts
  the fact.
\item[Templates] A table of the templates in every module. Create new
  Template opens an editor with a slot list and a live preview of the
  \fn{deftemplate} it will send; right-clicking a row removes a
  template.
\item[Functions] Every function the engine knows, with its usage text.
\item[Rete] The network viewer, described below.
\item[Log] Every message that passes through the message router, from
  every channel, with its time, channel and type: commands, engine
  output, results and errors. This is where to look when a command
  typed elsewhere seems to do nothing.
\item[Settings] Two panes. \emph{Engine} has check boxes for the profile
  and watch switches, each of which sends the corresponding command;
  \emph{Shell} sets the console font, size and colours. Save Changes
  writes the preferences.
\end{description}

The File menu has \emph{Batch File}, which opens a file chooser and
sends \fn{(batch path)} for the file you pick, \emph{Close Gui} and
\emph{Quit}.

\section{The network viewer}

The Rete tab, and the \fn{(view)} function typed at any prompt, draw the
whole RETE network as a graph. Object type nodes are orange, left input
adapters cyan, joins green, alpha nodes red and terminal nodes black;
lines out of a join are red, all others blue. The view zooms and pans, a
small radar panel shows which part of the network is visible, and
clicking a node prints its description (the constraint it tests or the
bindings it joins on, and its use count) in the text area below.
\emph{Reload View} redraws the network; \emph{Automatic Reload} redraws
it by itself whenever a rule is added or removed or working memory
changes, a few times a second at most.

Read the diagram top down: each template's node at the top, the alpha
tests of every rule that uses it below, the joins where patterns meet,
and one terminal per rule. Shared alpha nodes appear once with several
outgoing lines, which is the cheapest way to see whether two rules are
sharing work (Chapter~\ref{ch:analysis}).
